% Auxiliary register-restriction identities; verification marks withheld.

\subsection{Restrictions with excluded registers}

Let $P,Q$ be arbitrary sets of party labels. A register consists of
an owner and a complex inner product space. For an ordered finite
register list $a$, let $M(a)$ be its right-associated tensor memory,
including a terminal factor $\C$; the empty list has memory $\C$.
For a selection $f:Q\to\{0,1\}$, write $a_f$ for the ordered sublist
of registers whose owners are selected. An owner map $q:P\to Q$
changes the labels in $a$, retaining its spaces and order; denote
the resulting list by $q_*a$.

\begin{theorem}[An owner map with no selected image]
    \label{thm:peps_excluded_mapped_owners}
    \lean{TNLean.PEPS.PairEffect.Layout.restrict_mapOwner_eq_nil}
    Let $q:P\to Q$ be any map and $f:Q\to\{0,1\}$ any selection.
    Suppose $f(q(p))=0$ for every $p\in P$. Then, for every register
    list $a$ with owners in $P$, the selected sublist $(q_*a)_f$
    is empty. No injectivity, finiteness, or nonemptiness assumption
    on the party sets is required.
\end{theorem}

\begin{proof}
    If a register belonged to the selected sublist, it would come
    from a register with some original owner $p\in P$ and would
    satisfy $f(q(p))=1$. This contradicts the hypothesis.
\end{proof}

A source-free word is a finite composition of local linear operations,
register exchanges, and identities, containing no pair-source
preparations. Its restriction to selected parties retains their local
operations and registers and restricts the exchanges accordingly.
For a source-free word $W$ from register list $a$ to register list $b$,
write $W|_f:M(a_f)\to M(b_f)$ for its selected operator. If $r$ is a
spectator register list, write $\widetilde W$ for the word obtained
by adjoining $r$ unchanged to every step of $W$, from $ra$ to $rb$.
Concatenation of register lists is denoted by juxtaposition.

\begin{theorem}[Excluded spectator registers preserve the restriction]
    \label{thm:peps_excluded_spectator_restriction}
    \lean{TNLean.PEPS.PairEffect.Word.restrict_frameList_eval_heq_of_eq_nil}
    Let $f:Q\to\{0,1\}$ be any selection, let $r$ be a register
    list with $r_f=\varnothing$, and let $W:a\to b$ be any
    source-free word with owners in $Q$. Under the canonical
    identifications $(ra)_f=a_f$ and $(rb)_f=b_f$,
    $\widetilde W|_f=W|_f$. The retained input and output lists may
    be arbitrary. No contraction bound or nonzero-dimension
    assumption on the registers is imposed.
\end{theorem}

\begin{proof}
    Induct on the spectator list $r$. For the empty list the two
    words coincide. For a nonempty list, its first register is
    excluded because $r_f=\varnothing$. Restricting the word framed
    by this register gives the same selected operator as restricting
    the remaining framed word. Apply the induction hypothesis to
    the remaining spectator registers.
\end{proof}

These identities concern the register selections used in
\cite[proof of Theorem~5.2]{OpenAI2026PolynomialPEPS},
\texttt{04-compression.tex}, lines 233--267 and 351--381.
